\documentclass[11pt,a4paper]{article}
\usepackage{../../common/preamble}

\begin{document}

\begin{center}
    {\LARGE \textbf{Передача данных. Помехоустойчивые коды}}\\[0.2cm]
    \rule{\textwidth}{1pt}
\end{center}

\section{Математическая модель канала передачи данных}

\subsection{Скорость передачи}

\begin{definitionbox}{Скорость передачи и пропускная способность}
\begin{enumerate}
    \item \textbf{Скорость передачи данных} --- количество информации, которое передается по каналу связи за единицу времени. Основной единицей измерения является бит в секунду.
    \item \textbf{Пропускная способность канала связи} --- наибольшая возможная скорость передачи данных, которую принципиально невозможно превысить.
\end{enumerate}
\end{definitionbox}

Процесс передачи дискретных данных во времени описывается базовым кинематическим соотношением:
$$I = v \cdot t$$
где $v$ --- средняя скорость передачи информации (называемая битрейтом), $t$ --- время передачи, $I$ --- информационный объем.

Естественно, в ходе развития каналов связи появилась необходимость в использовании единиц измерения скорости передачи данных, кратных битам в секунду. При этом важно отметитть, что для их обозначения используются десятичные приставки, а не двоичные, как, например, в случае измерения количества информации. Так, используются в основном килобиты в секунду ($1 \text{ кбит/с} = 1000 \text{ бит/c}$), мегабиты в секунду ($1 \text{ Мбит/с} = 10^6 \text{ бит/с}$) и гигабиты в секунду ($1 \text{ Гбит/с} = 10^9 \text{ бит/с}$), а также иногда бывают полезны измерения в байтах в секунду ($1 \text{ байт/с} = 8 \text{ бит/с}$) и мегабайтах в секунду ($1 \text{ Мбайт/с} = 10^6 \text{ байт/с}$).

\section{Метрика Хэмминга и обнаружение ошибок}

\subsection{Пространство двоичных слов и расстояние Хэмминга}
Рассмотрим множество всех цепочек длины $n$, состоящих из нулей и единиц. Также скажем, что их можно складывать следующим образом: в цепочке, являющейся результатом сложения двух цепочек, стоят нули в местах, где соответствующие элементы исходных цепочек совпадают, и единицы в остальных местах. Давайте формализуем это определение.

Пусть $V_n = \mathbb{F}_2^n$ --- векторное пространство двоичных последовательностей длины $n$ над полем из двух элементов $\mathbb{F}_2 = \{0, 1\}$. Операция сложения векторов $\vec{x}, \vec{y} \in \mathbb{F}_2^n$ определяется как покомпонентное сложение по модулю $2$ (XOR, обозначается $\oplus$):
$$\vec{x} \oplus \vec{y} = (x_1 \oplus y_1, \, x_2 \oplus y_2, \, \dots, \, x_n \oplus y_n)$$

\begin{definitionbox}{Вес и расстояние Хэмминга}
\begin{enumerate}
    \item \textbf{Весом Хэмминга} $w(\vec{x})$ вектора $\vec{x}$ называется число его ненулевых компонент (количество единиц).
    \item \textbf{Расстоянием Хэмминга} $d(\vec{x}, \vec{y})$ между двумя кодовыми словами одинаковой длины называется число позиций, в которых эти векторы различаются:
    $$d(\vec{x}, \vec{y}) \defeq w(\vec{x} \oplus \vec{y}) = \sum_{i=1}^n (x_i \oplus y_i)$$
\end{enumerate}
\end{definitionbox}

Функция $d(\vec{x}, \vec{y})$ является \textbf{метрикой} (она же - расстояние) в топологическом смысле, так как удовлетворяет трем аксиомам:
\begin{enumerate}
    \item Аксиома тождества: $d(\vec{x}, \vec{y}) \ge 0$, причем $d(\vec{x}, \vec{y}) = 0 \iff \vec{x} = \vec{y}$.
    \item Аксиома симметрии: $d(\vec{x}, \vec{y}) = d(\vec{y}, \vec{x})$.
    \item Неравенство треугольника: $\forall \vec{x}, \vec{y}, \vec{z} \in \mathbb{F}_2^n \implies d(\vec{x}, \vec{z}) \le d(\vec{x}, \vec{y}) + d(\vec{y}, \vec{z})$.
\end{enumerate}

\subsection{Кодовое расстояние и корректирующая способность кодов}
Пусть $C \subset \mathbb{F}_2^n$ --- множество разрешенных кодовых слов. 
\begin{definitionbox}{Минимальное кодовое расстояние}
Минимальным расстоянием Хэмминга кода $C$ называется величина:
$$d_{\min} \defeq \min_{\substack{\vec{x}, \vec{y} \in C \\ \vec{x} \neq \vec{y}}} d(\vec{x}, \vec{y})$$
\end{definitionbox}

\begin{theorembox}
\textbf{Теорема о корректирующей способности кода:}
\begin{enumerate}
    \item Код $C$ гарантированно \textbf{обнаруживает} любые конфигурации до $s$ ошибок включительно тогда и только тогда, когда:
    $$d_{\min} \ge s + 1$$
    \item Код $C$ гарантированно \textbf{исправляет} любые конфигурации до $r$ ошибок включительно (декодирование по принципу максимального правдоподобия / минимального расстояния) тогда и только тогда, когда:
    $$d_{\min} \ge 2r + 1$$
\end{enumerate}
\end{theorembox}

\textit{Геометрическая интерпретация:} Ошибка кратности $r$ смещает переданное кодовое слово $\vec{x}$ в точку $\vec{y}$ внутри гипершара радиуса $r$: $B_r(\vec{x}) = \{\vec{z} \mid d(\vec{x}, \vec{z}) \le r\}$ (по сути своей такой шар является множеством всех слов, которые отличаются от заданного не более чем в $r$ позициях). Чтобы при декодировании не возникало неоднозначности, шары радиуса $r$ вокруг любых двух разрешенных слов $\vec{x}_1, \vec{x}_2$ не должны пересекаться (иначе мы просто не сможем определить, в центр какого шара смещать слово):
$$B_r(\vec{x}_1) \cap B_r(\vec{x}_2) = \varnothing \iff d(\vec{x}_1, \vec{x}_2) \ge 2r + 1$$

\begin{center}
\begin{tikzpicture}[scale=0.9]
    \draw[thick, fill=black!5] (0, 0) circle (1.4);
    \filldraw[black] (0, 0) circle (2pt) node[below=2pt] {$\vec{x}_1 \in C$};
    \draw[dashed, ->] (0, 0) -- (0.99, 0.99) node[midway, above left] {$r$};
    
    \draw[thick, fill=black!5] (4, 0) circle (1.4);
    \filldraw[black] (4, 0) circle (2pt) node[below=2pt] {$\vec{x}_2 \in C$};
    \draw[dashed, ->] (4, 0) -- (4.99, 0.99) node[midway, above left] {$r$};
    
    \draw[thick, Stealth-Stealth] (0, -0.6) -- (4, -0.6) node[midway, below] {$d_{\min} \ge 2r + 1$};
    
    \filldraw[black] (1.1, 0.3) circle (1.5pt) node[above right] {$\vec{y}$};
    \draw[dotted, thick, ->] (0, 0) -- (1.1, 0.3);
    \node[right] at (5.7, 0.5) {\footnotesize $\vec{y}$ ближе к $\vec{x}_1$, чем к $\vec{x}_2 \implies$};
    \node[right] at (5.7, -0.1) {\footnotesize декодируется однозначно в $\vec{x}_1$};
\end{tikzpicture}
\end{center}

\subsection{Простейшие помехоустойчивые коды}
\begin{enumerate}
    \item \textbf{Бит четности (Parity Bit):} К блоку из $k$ бит добавляется 1 контрольный бит $c = b_1 \oplus b_2 \oplus \dots \oplus b_k$. Сумма всех бит кодового слова по модулю 2 всегда равна нулю. 
    $$d_{\min} = 2 \implies s = 1, \; r = 0$$
    Код обнаруживает одиночные (и любые нечетные) ошибки, но \textbf{не способен исправить} ни одной ошибки и не замечает четное число искажений.
    \item \textbf{Код с троекратным повторением (Repetition Code):} Бит $0 \to 000$, $1 \to 111$. Разрешенные слова $\{000, 111\}$. 
    $$d_{\min} = 3 \implies s = 2, \; r = 1$$
    Позволяет исправлять $1$ ошибку методом мажоритарного голосования (например, $001 \to 000$).
\end{enumerate}

\section{Линейный код Хэмминга $(7, 4)$}

\subsection{Постановка задачи}
Как построить код с $d_{\min} = 3$ (исправление $r = 1$ ошибки), расходуя на избыточность минимум бит? 

В коде Хэмминга $(7, 4)$:
\begin{itemize}
    \item Общая длина кодового слова: $n = 7$.
    \item Число информационных бит: $k = 4$.
    \item Число контрольных бит: $m = n - k = 3$.
\end{itemize}

\subsection{Алгоритм формирования кода}
Номера позиций нумеруются слева направо числами от $1$ до $7$. 
Правило:
\begin{enumerate}
    \item Позиции, номера которых являются степенями двойки, отводятся под контрольные биты: $1, 2, 4$ ($c_1, c_2, c_4$).
    \item Оставшиеся позиции заполняются битами данных: $3, 5, 6, 7$ ($d_1, d_2, d_3, d_4$).
\end{enumerate}

\begin{table}[h!]
\centering
\begin{tabular}{@{}lccccccc@{}}
\toprule
\textbf{Номер позиции ($i$)} & \textbf{1} & \textbf{2} & \textbf{3} & \textbf{4} & \textbf{5} & \textbf{6} & \textbf{7} \\ \midrule
\textbf{Двоичный индекс ($i_2$)} & $001_2$ & $010_2$ & $011_2$ & $100_2$ & $101_2$ & $110_2$ & $111_2$ \\
\textbf{Назначение бита} & $c_1$ & $c_2$ & $d_1$ & $c_4$ & $d_2$ & $d_3$ & $d_4$ \\ \bottomrule
\end{tabular}
\end{table}

Каждый информационный индекс раскладывается в сумму степеней двойки (притом, как можно догадаться из возможности перевода чисел между двоичной и десятичной системами счисления, единственным образом):
$$3 = 1 + 2, \quad 5 = 1 + 4, \quad 6 = 2 + 4, \quad 7 = 1 + 2 + 4$$

Контрольный бит с номером $2^j$ осуществляет проверку на четность всех позиций, в двоичной записи номеров которых на $j$-м месте стоит единица:
\begin{align*}
    c_1 &= d_1 \oplus d_2 \oplus d_4 \\
    c_2 &= d_1 \oplus d_3 \oplus d_4 \\
    c_4 &= d_2 \oplus d_3 \oplus d_4
\end{align*}

\subsection{Исправление ошибки}
Пусть на приемной стороне получен вектор $\vec{y} = (y_1, y_2, y_3, y_4, y_5, y_6, y_7)$. Вычисляются три бита \textbf{синдрома ошибки} $\vec{s} = (s_4, s_2, s_1)_2$:
\begin{align*}
    s_1 &= y_1 \oplus y_3 \oplus y_5 \oplus y_7 \\
    s_2 &= y_2 \oplus y_3 \oplus y_6 \oplus y_7 \\
    s_4 &= y_4 \oplus y_5 \oplus y_6 \oplus y_7
\end{align*}

\begin{definitionbox}{Вектор синдрома}
Число $S \defeq s_4 \cdot 2^2 + s_2 \cdot 2^1 + s_1 \cdot 2^0$ указывает состояние канала:
\begin{itemize}
    \item Если $S = 0$ ($000_2$) --- ошибок не обнаружено (или их кратность $> 2$).
    \item Если $S \in \{1, 2, \dots, 7\}$ --- значение $S$ \textbf{в точности равно номеру ошибочного бита}. Для исправления достаточно инвертировать бит $y_S \coloneqq y_S \oplus 1$.
\end{itemize}
\end{definitionbox}

\section{Практика}

\subsection{ЕГЭ №7: Архивирование и передача данных}
\begin{taskbox}{}
Музыкальный фрагмент был оцифрован и записан в виде файла без использования сжатия данных. Получившийся файл был передан в город А по каналу связи за 15 секунд. Затем тот же музыкальный фрагмент был оцифрован повторно с разрешением в 2 раза выше и частотой дискретизации в 1,5 раза меньше, чем в первый раз. Сжатие данных не производилось. Полученный файл был передан в город Б; пропускная способность канала связи с городом Б в 2 раза выше, чем канала связи с городом А. Сколько секунд длилась передача файла в город Б?

В ответе запишите только целое число, единицу измерения писать не нужно.
\end{taskbox}

\textbf{Решение:}
\begin{enumerate}
    \item Формулки...
    \begin{itemize}
        \item Информационный объем несжатого аудиофайла:
        $$I = k \cdot f \cdot i \cdot T$$
        где $k$ --- количество каналов, $f$ --- частота дискретизации, $i$ --- разрешение, $T$ --- длительность.
        \item Время передачи файла по каналу с пропускной способностью $v$:
        $$t = \frac{I}{v} = \frac{k \cdot f \cdot i \cdot T}{v}$$
    \end{itemize}

    \item
    Пусть величины с индексом $1$ относятся к передаче в город А, а с индексом $2$ --- к передаче в город Б.
    По условию задачи:
    \begin{itemize}
        \item Музыкальный фрагмент и число каналов не менялись: $T_2 = T_1$, $k_2 = k_1$.
        \item Разрешение увеличено в 2 раза: $i_2 = 2 \cdot i_1$.
        \item Частота дискретизации уменьшена в 1.5 раза: $f_2 = \frac{f_1}{1.5} = \frac{2}{3} f_1$.
        \item Скорость канала связи увеличилась в 2 раза: $v_2 = 2 \cdot v_1$.
    \end{itemize}

    \item Выразим время передачи второго файла $t_2$ через время передачи первого $t_1 = 15\text{ с}$:
    $$\frac{t_2}{t_1} = \frac{I_2 / v_2}{I_1 / v_1} = \frac{I_2}{I_1} \cdot \frac{v_1}{v_2}$$
    
    Найдем относительное изменение объема файла:
    $$\frac{I_2}{I_1} = \frac{f_2}{f_1} \cdot \frac{i_2}{i_1} = \frac{1}{1.5} \cdot 2 = \frac{2}{3} \cdot 2 = \frac{4}{3}$$

    Учтем изменение пропускной способности канала:
    $$\frac{t_2}{t_1} = \frac{I_2}{I_1} \cdot \frac{v_1}{v_2} = \frac{4}{3} \cdot \frac{1}{2} = \frac{2}{3}$$

    \item Таким образом, получаем время:
    $$t_2 = t_1 \cdot \frac{2}{3} = 15 \cdot \frac{2}{3} = 10 \text{ с}$$
\end{enumerate}
\textbf{Ответ:} 10

\subsection{Просто прикольная задача на коды Хэмминга}
\begin{taskbox}{Поиск и исправление ошибки}
По каналу связи с помехами передавалось 4 бита данных, закодированных кодом Хэмминга $(7, 4)$. На приемной стороне была получена двоичная последовательность:
$$\vec{y} = 0111110_2$$
Известно, что в канале произошло не более одной ошибки. 
\begin{enumerate}
    \item Определите, произошла ли ошибка при передаче. Если да, найдите номер ошибочной позиции.
    \item Восстановите исходные 4 бита данных, которые отправлял передатчик.
\end{enumerate}
\end{taskbox}

\textbf{Решение:}
\begin{enumerate}
    \item Выпишем полученные биты по их позициям $1..7$:
    $$y_1 = 0, \quad y_2 = 1, \quad y_3 = 1, \quad y_4 = 1, \quad y_5 = 1, \quad y_6 = 1, \quad y_7 = 0$$
    \item Вычислим компоненты проверочного синдрома $\vec{s} = (s_4, s_2, s_1)$:
    \begin{align*}
        s_1 &= y_1 \oplus y_3 \oplus y_5 \oplus y_7 = 0 \oplus 1 \oplus 1 \oplus 0 = 0 \\
        s_2 &= y_2 \oplus y_3 \oplus y_6 \oplus y_7 = 1 \oplus 1 \oplus 1 \oplus 0 = 1 \\
        s_4 &= y_4 \oplus y_5 \oplus y_6 \oplus y_7 = 1 \oplus 1 \oplus 1 \oplus 0 = 1
    \end{align*}
    \item Собираем номер ошибочной позиции в двоичной системе счисления:
    $$S = (s_4 s_2 s_1)_2 = 110_2 = 1 \cdot 4 + 1 \cdot 2 + 0 \cdot 1 = 6$$
    Это число отлично от нуля, следовательно, ошибка есть и произошла она в бите $y_6$.
    \item Инвертируем ошибочный бит для восстановления правильного кодового слова:
    $$\vec{x}_{\text{испр}} = 01111\mathbf{0}0_2$$
    \item Извлекаем исходные информационные биты (позиции $3, 5, 6, 7$):
    $$d_1 = x_3 = 1, \quad d_2 = x_5 = 1, \quad d_3 = x_6 = 0, \quad d_4 = x_7 = 0$$
    Исходное сообщение: $1100_2$.
\end{enumerate}
\textbf{Ответ:} Ошибка в 6-м бите; исходные данные: $1100$.

\end{document}